% This file is intended to be included via \input{} or \include{}.
% It contains no preamble and no document environment.

\section*{Appendix A --- Full Poster-Style Walkthrough (Dual-Flow, Detail-Preserving)}

This is a detail-preserving reorganization of the earlier full poster style Appendix companion.\par

We keep the rich blocks, callouts, and motivations from the previous version.\par

We reorganize them into the dual-flow learning sequence from the Math Primer concept map.\par

No mathematics is changed; only the teaching order is changed.\par

Two flows:\par

Content flow: explicit objects, reductions, category E\par

Structure flow: TDG/shapes, navigation, category S/N These synchronize at the fibration and proceed to the combined system, semantics, and theorems.\par

\section*{[PHASE] PHASE I --- EXPLICIT CONTENT (WHAT THE AI CAN DO)}

\subsection*{[PHASE] BLOCK C1 --- Tacit vs Explicit Worlds}

Motivation: The AI must operate entirely in an inspectable space; tacit human reasoning cannot be accessed or simulated.\par

Tacit World (Human Only): intuition, unstated assumptions, judgment.\par

Explicit World (Human + AI): formal statements, symbolic expressions, explicit computations.\par

Note: Callout: FRFP begins by drawing a hard boundary: AI stays in the explicit world.\par

\subsection*{[PHASE] BLOCK C2 --- RI, EC, ED: Atomic Explicit Objects}

Motivation: We need a controlled vocabulary for every explicit move an AI makes.\par

RI: user/human-supplied explicit representations (definitions, assumptions, givens)\par

EC: explicit computations (symbolic, algebraic, numeric transforms)\par

ED: explicit diagnostics (clarifications, warnings, notes that do not change math)\par

 Interpretation: These are the atoms of explicit reasoning.\par

\subsection*{[PHASE] BLOCK C3 --- Pipelines: Composite Explicit Objects}

Motivation: Reasoning is not atomic; it is organized into workflows.\par

A pipeline is an explicit workflow built from RI/EC/ED.\par

Pipelines are inspectable and shareable.\par

Pipelines are the primary objects the AI transforms.\par

Key point: Key Point: If the AI reasons, it does so via a pipeline.\par

\subsection*{[PHASE] BLOCK C4 --- Explicit Reduction ($\to\_\{e\}$)}

Motivation: The explicit world needs a consistent normalization system.\par

Explicit reduction:\par

simplifies/normalizes explicit content\par

distributes through workflow structure (later made precise with TDG)\par

is designed to terminate and be locally confluent\par

 Result: Explicit workflows always normalize (to stable normal forms).\par

\subsection*{[PHASE] BLOCK C5 --- Explicit Category E (Content Category)}

Motivation: By treating pipelines as categorical objects, we gain mathematical control over workflow behavior.\par

Objects: explicit pipelines (composite explicit objects)\par

Morphisms: explicit reductions + structural identities\par

Note: Clarification: Atomic objects (RI/EC/ED) admit only identity morphisms; the nontrivial morphisms act on composite pipelines.\par

Key point: Key Benefit: Composition models multi-step reasoning; identities model do nothing. \par

\subsection*{[PHASE] BLOCK C6 --- Why This Algebraic Structure Enables Proofs}

Motivation: Proofs require stable rules for how workflows transform.\par

Because E has lawful composition and identity and a reduction relation:\par

we can state and prove normalization and confluence results\par

we can compare workflows via normal forms\par

we can preserve meaning under explicit transformations\par

Note: Callout: This is what algebraic structure suitable for formal reasoning means in FRFP: pipelines behave like mathematical objects with rule-governed transformations.\par

\section*{[PHASE] PHASE II --- PURE STRUCTURE (FORM WITHOUT CONTENT)}

\subsection*{[PHASE] BLOCK S1 --- TDG: The Structural Grammar}

Motivation: We need a grammar that guarantees workflows are well-structured and manipulable.\par

TDG provides constructors (workflow form):\par

chain(Q , Q , ) -- sequential steps\par

branch(Q , Q , ) -- case splits / alternatives\par

parallel(Q , Q , ) -- independent subwork\par

loop(Q\_cond, Q\_body) -- iterative refinement\par

RI, EC, ED -- atomic leaves\par

 Interpretation: TDG is the syntax engine for building structured pipelines.\par

\subsection*{[PHASE] BLOCK S2 --- Shapes: The Structural Skeleton}

Motivation: Separate structure from content so we can reorganize workflows without altering the math.\par

A shape is a pipeline s structure with payload removed.\par

Examples:\par

chain\par

branch(chain, chain)\par

loop(chain, branch(chain, chain))\par

 Poster Metaphor: Shapes are folders (structure); pipelines are files (content).\par

\subsection*{[PHASE] BLOCK S3 --- Shape Category S}

Motivation: To reason about structure formally, shapes must live in a category.\par

Objects: shapes\par

Morphisms: structure-preserving maps (identities + composition)\par

 Outcome: Structure becomes a first-class mathematical object.\par

\section*{[PHASE] PHASE III --- NAVIGATION (STRUCTURE-ONLY CHANGE)}

\subsection*{[PHASE] BLOCK N1 --- Navigation Category and Groupoid Intuition}

Motivation: Humans frequently ask to refactor workflows; we need a safe, reversible structural editing language.\par

Navigation consists of invertible structural rewrites (examples):\par

reorder branches\par

flatten nested branches\par

regroup parallels\par

distribute loops\par

 Key Idea: Navigation changes shape only, not explicit content.\par

Note: Why groupoid? Every navigation step has an inverse (can be undone).\par

\subsection*{[PHASE] BLOCK N2 --- Why Navigation Is Not a Group}

Motivation: Navigation operations are invertible, so why not model them as a group?\par

A group describes symmetries of a single object. Navigation, however, describes invertible rewrites between many different shapes.\par

Shapes are distinct objects\par

Navigation morphisms go between shapes\par

Each rewrite has an inverse, but not a single global carrier\par

Note: Formal insight: Navigation forms a groupoid (a category where every morphism is invertible), not a group.\par

 Callout (one-sentence formal clarification): Navigation is modeled as a groupoid rather than a group because structural rewrites are invertible but act between many distinct shapes, not as symmetries of a single shape.\par

Note: Clarification --- Why groupoid and not group: A group describes invertible operations on one fixed object, whereas navigation must describe invertible rewrites between many different shapes; this multiplicity of objects forces a groupoid rather than a group.\par

\section*{[PHASE] PHASE IV --- SYNCHRONIZATION (WHERE CONTENT MEETS STRUCTURE)}

\subsection*{[PHASE] BLOCK F1 --- Fibers: All Pipelines Inside a Folder }

Motivation: Grouping pipelines by structure ensures restructuring operations know exactly where each pipeline belongs.\par

For a shape S:\par

fiber(S) = \{ e E | shape(e) = S \}\par

 Metaphor: The fiber over S is the folder contents---all pipelines with structure S.\par

Note: Callout: e fiber(shape(e)) means: pipeline e lives in the folder labeled by its shape.\par

\subsection*{[PHASE] BLOCK F2 --- Shape Fibration p : E $\to$ S}

Motivation: We need a canonical rule that assigns each pipeline to its structural folder and keeps them aligned.\par

The fibration maps each explicit pipeline to its shape:\par

p(e) = shape(e)\par

 Key Idea: Fibration is the alignment mechanism between structure and content.\par

\subsection*{[PHASE] BLOCK F3 --- Cartesian Lifting (Safe Transport)}

Motivation: When a folder (shape) moves, every file (pipeline) must move safely and canonically.\par

Given navigation $\sigma$ : S $\to$ S' and e with p(e)=S, cartesian lifting produces e' with p(e')=S':\par

 e --------$\to$ e'
 | |
 p p 
 S ---$\sigma$$\to$ S'\par

 Callout: Liftings ensure structural refactoring never corrupts explicit math.\par

 Core Safety Insight: Cartesian lifting plays the same safety role for structure/content separation that the tacit/explicit boundary plays for human/AI separation.\par

\subsection*{[PHASE] BLOCK F4 --- Navigation Action on Pipelines}

Motivation: We need one operation that applies structural rewrites directly to explicit workflows.\par

Define:\par

$\sigma \cdot e$ = e'\par

where e' is the cartesian lift of $\sigma$ applied to e.\par

Note: Clarification: Cartesian lifting states that such an e' exists and is unique; the navigation action is the explicit operation that takes e to that e'.\par

Key point: Outcome: Structure changes; content is preserved.\par

\section*{[PHASE] PHASE V --- THE COMBINED SYSTEM (STRUCTURE + CONTENT)}

\subsection*{[PHASE] BLOCK G1 --- Why Combine Navigation and Explicit Computation?}

Motivation: Real reasoning evolves both structure and content; we need one framework that tracks both.\par

Navigation handles form, explicit reduction handles math. The combined system tracks both coherently.\par

\subsection*{[PHASE] BLOCK G2 --- Grothendieck Construction / Semidirect Product N $\ltimes$ E}

Motivation: We need a single system where each step records (i) structure change and (ii) the corresponding content-following change.\par

In N $\ltimes$ E:\par

Objects: (S, e) where e lies over S\par

Morphisms: ($\sigma$, f) where $\sigma$ changes shape and f is the induced/compatible content morphism\par

Note: Key Idea: A morphism remembers both what structure changed and how content followed. \par

\subsection*{[PHASE] BLOCK G3 --- Composition and Combined Rewrite System}

Motivation: The system must remain predictable when steps are chained.\par

Composition:\par

shape part composes ($\sigma$ $\sigma$ )\par

content part composes after reindexing / lifting\par

Combined rewrite:\par

explicit reduction steps\par

navigation steps\par

mixed sequences\par

 Outcome: Mixed workflow evolution is well-defined.\par

\subsection*{[PHASE] BLOCK G4 --- Confluence (Order Doesn t Matter)}

Motivation: Humans and AI interleave edits in varying orders; results should not depend on interaction order.\par

Assumptions ensure:\par

mixed sequences converge\par

normal forms are stable (up to navigation isomorphism)\par

Key point: Practical Benefit: You can restructure while computing and still get consistent outcomes.\par

\section*{[PHASE] PHASE VI --- SEMANTICS \& TACIT CORRECTNESS}

\subsection*{[PHASE] BLOCK M1 --- Why Semantics Is Needed}

Motivation: Explicit pipelines must ultimately produce meaning interpretable by humans.\par

Semantics is the bridge from explicit computation to human-level meaning.\par

\subsection*{[PHASE] BLOCK M2 --- Semantic Interpretation \(\llbracket$ \rrbracket\)}

Motivation: We need a rule that maps an explicit pipeline to its outcome.\par

\(\llbracket$ \rrbracket\) : E $\to$ Outcome\par

Note: Callout: This explains what a pipeline means, not just what it does.\par

\subsection*{[PHASE] BLOCK M3 --- TE and HFD (Human Roles)}

Motivation: Some judgments cannot be mechanized; humans must evaluate and decide.\par

TE: tacit evaluation\par

HFD: human final decision\par

These occur outside the explicit world.\par

\subsection*{[PHASE] BLOCK M4 --- RB Functor (Boundary / Gateway)}

Motivation: We need a controlled handoff from explicit to tacit interpretation.\par

RB is the explicit-to-tacit boundary mechanism:\par

prepares outputs for human interpretation\par

prevents hidden computation at the boundary\par

 Metaphor: RB is the airlock between explicit and tacit worlds.\par

\subsection*{[PHASE] BLOCK M5 --- Navigation Invariance of Meaning}

Motivation: Refactoring structure must not change what the result means.\par

If e and e' differ only by navigation action:\par

$\sigma \cdot e$ $\cong$ e'\par

then:\par

\(\llbracket$e\rrbracket\) = \(\llbracket$e'\rrbracket\)\par

 Outcome: Structural editing is meaning-safe.\par

\subsection*{[PHASE] BLOCK M6 --- Tacit Correctness}

Motivation: Correctness ultimately lives in human interpretation.\par

A pipeline is tacit-correct when:\par

explicit normalization behaves well\par

semantic meaning is well-defined\par

TE and HFD validate it\par

Key point: Result: FRFP outputs are human-verifiable.\par

\section*{[PHASE] PHASE VII --- FRFP THEOREMS (WHAT WE CAN NOW PROVE)}

\subsection*{[PHASE] BLOCK T1 --- Main Theorems (Poster Summary)}

Motivation: The architecture exists to make strong guarantees.\par

Theorems establish:\par

Well-formedness (pipelines stay valid)\par

Explicit completeness (AI actions live in E)\par

Minimality / initiality (the explicit world is canonical)\par

Canonicality (stable normal forms)\par

Confluence (order independence)\par

Decidability (computability of reductions)\par

Correctness preservation (meaning preserved)\par

Semantic/tacit correctness (human interpretability)\par

 Outcome: FRFP supports safe, inspectable, refactorable human--AI collaboration.\par

\section*{[PHASE] PHASE VIII --- EXAMPLES \& META-REMARKS}

\subsection*{[PHASE] BLOCK X1 --- Worked Examples (Why They Matter)}

Motivation: Examples demonstrate that the machinery is usable, not just formal.\par

Examples typically show:\par

chain pipelines reducing cleanly\par

branching workflows reorganized by navigation\par

loop/branch interaction and safe liftings\par

mixed navigation + explicit reduction reaching the same normal form\par

semantics producing human-interpretable outcomes\par

\subsection*{[PHASE] BLOCK X2 --- Meta-Remarks and Implementation Notes}

Motivation: Implementation needs guidance, not just proofs.\par

Key implementation principles:\par

represent pipelines and shapes as typed, inspectable data\par

separate navigation (structure) from reduction (content)\par

enforce explicit-only computation\par

surface semantics only at the human boundary\par

\section*{FINAL BIG PICTURE (DUAL-FLOW)}

Develop explicit content and pure structure independently, synchronize them via fibration and cartesian lifting, combine them into one coherent system, interpret outcomes semantically for humans, and prove correctness and stability by construction.\par

End of Detail-Preserving Restructured Appendix A\par
